\chapter{Závěr}
\label{ch:conclusion}

Tato diplomová práce se zabývala návrhem, implementací a ověřením modulu pro automatickou analýzu studentských zdrojových kódů pomocí velkých jazykových modelů v systému Kelvin.
Cílem bylo poskytnout vyučujícím podpůrný nástroj, který jim pomůže rychleji se zorientovat v odevzdaných řešeních a upozorní na problematická místa v kódu, aniž by nahrazoval jejich odborný úsudek.

Na základě analýzy dostupných modelů, způsobů nasazení a promptovacích strategií bylo zvoleno on-premise nasazení modelu Qwen3-Coder:30B prostřednictvím nástroje Ollama, který zajišťuje ochranu studentských dat a nulové marginální náklady na zpracování odevzdání.
Implementace zahrnuje backendový Django modul provádějící asynchronní LLM analýzu souběžně s evaluační pipeline a frontendové Vue.js komponenty, které učiteli zobrazuje navrhované komentáře přímo u zdrojového kódu s možností jejich přijetí, úpravy nebo zamítnutí, či umožňuje upravovat celkové prompty používané pro generování komentářů.
Studentovi se zobrazují pouze komentáře explicitně schválené učitelem, čímž je zajištěna plná kontrola pedagoga nad publikovanou zpětnou vazbou.

Experimentální měření potvrdilo, že chain-of-thought přístup se třemi kroky analýzy produkuje pečlivěji ověřené nálezy díky filtraci falešně pozitivních zjištění, zatímco one-shot přístup je přibližně 8,5krát rychlejší a vhodnější pro cloudové API z hlediska nákladů.
Zpětná vazba získaná od učitelů prostřednictvím samostatné evaluační aplikace na reálných studentských datech z předmětu UPR potvrdila schopnost zvoleného modelu identifikovat relevantní problémy ve studentských řešeních.
Hlavním omezením zůstává pravděpodobnostní povaha jazykových modelů, které neposkytují formální záruku správnosti a jejichž výstupy mohou obsahovat nepřesnosti.

Z pohledu dalšího rozvoje se nabízí především začlenění výsledků evaluační pipeline do vstupu LLM analýzy, které by modelu umožnilo pracovat i s výstupy kompilátoru, a využití shromážděných hodnocení od učitelů pro fine-tuning modelu za účelem zvýšení relevance generovaných komentářů.
Přínosné by rovněž bylo rozsáhlejší pilotní nasazení v produkčním provozu a dlouhodobější vyhodnocení skutečného dopadu na efektivitu hodnocení a kvalitu zpětné vazby poskytované studentům.

\endinput